% Supplementary Material for "Task-Oriented Semantic Granularity Selection for
% Bandwidth-Constrained V2V Cooperative Perception".
% Standalone: compiles on its own (tectonic supplementary.tex). Every number here is the same
% number the repository products carry; nothing was re-derived for this file.
\documentclass[journal]{IEEEtran}
\usepackage{amsmath,amssymb}
\usepackage{graphicx}
\usepackage{booktabs}
\usepackage{multirow}
\usepackage{cite}
\usepackage{xcolor}
\usepackage{url}
\graphicspath{{figures/}}
\newcommand{\method}{CA-TOSG}

\begin{document}
\title{Supplementary Material\\ Task-Oriented Semantic Granularity Selection for
Bandwidth-Constrained V2V Cooperative Perception}
\maketitle

\noindent This document carries material moved out of the main paper for length. It is
\emph{not} a second set of results: the two appendices below are the prior-protocol arms
labelled as such in the main text, and the cue table is the definition list of the 21
committed ego-side cues. Section and label names are preserved so that the main text's
pointers resolve against this document.

\section{Ego-Side Online Task Cues: full definition list}\label{sec:cues_full}
The selector's task-cue vector $z_t\in\mathbb{R}^{21}$ is defined in full below; the main
paper lists five representative cues and points here.

\begin{table}[t]
\centering
\footnotesize
\setlength{\tabcolsep}{3pt}
\caption{Ego-side online task cues used by the \method{} selector ($d=21$). All cues are computed locally on the ego vehicle before message selection, removing the causality issue that would arise from depending on the collaborator's state.}
\label{tab:cues}
\begin{tabular}{lll}
\toprule
Category & Cue & Definition \\
\midrule
\multirow{2}{*}{Agent context} & \texttt{num\_cavs} & number of CAVs in scene \\
                               & \texttt{ego\_num\_objects} & ego-detected objects \\
\midrule
\multirow{2}{*}{Frame shape}   & \texttt{ego\_origin\_lidar\_shape\_0/1} & raw LiDAR tensor shape \\
\midrule
\multirow{4}{*}{Range stats}   & \texttt{pcd\_num\_points} & total point count \\
                               & \texttt{pcd\_mean\_range}    & mean LiDAR range \\
                               & \texttt{pcd\_max\_range}     & max LiDAR range \\
                               & \texttt{pcd\_std\_range}     & range std.\ dev. \\
\midrule
\multirow{4}{*}{Distance bins} & \texttt{pcd\_near\_20m}      & points within 20 m \\
                               & \texttt{pcd\_mid\_20\_50m}   & points 20--50 m \\
                               & \texttt{pcd\_far\_50\_80m}   & points 50--80 m \\
                               & \texttt{pcd\_very\_far\_80m} & points beyond 80 m \\
\midrule
\multirow{4}{*}{Quadrant cnt.} & \texttt{pcd\_front\_points}  & front-half points \\
                               & \texttt{pcd\_back\_points}   & rear-half points \\
                               & \texttt{pcd\_left\_points}   & left-half points \\
                               & \texttt{pcd\_right\_points}  & right-half points \\
\midrule
\multirow{2}{*}{Long-range}    & \texttt{pcd\_front\_far\_30m} & front points $>30$ m \\
                               & \texttt{pcd\_front\_far\_50m} & front points $>50$ m \\
\midrule
\multirow{3}{*}{Density bins}  & \texttt{pcd\_density\_0\_20}    & density 0--20 m \\
                               & \texttt{pcd\_density\_20\_50}   & density 20--50 m \\
                               & \texttt{pcd\_density\_50\_80}   & density 50--80 m \\
\bottomrule
\end{tabular}
\end{table}


\section{When Are Perception Cues Necessary? LDPC Cliff versus JSCC Graceful Degradation}\label{sec:jscc_aware}

\noindent\emph{Protocol note.} A \textbf{prior-protocol arm}, \emph{not} re-evaluated under the frozen P2 protocol: its selector is trained and evaluated inside the codec-comparison pipeline, sharing the retired engine's BLER lookup and $200$-realisation CSI convention. Its numbers may not appear in the same sentence as a frozen-replay number, and no main-text claim rests on them.

Under LDPC + QAM a one-line threshold tracks the learned selector (Section~\ref{sec:threshold}); does that equivalence survive a feature codec \emph{without} a sharp SNR cliff? We replace the feature branch with importance-map JSCC~\cite{sheng2024importance} and re-run the per-frame pipeline: per-frame JSCC F1 is re-cached at each SNR ($1{,}000$ frames per SNR, AWGN and Rayleigh), the $\{L,\text{Feature-JSCC}\}$ oracle relabelled and the selector retrained---no perception-backbone change.

The mechanism is a codec property. The LDPC + QAM feature F1 cliffs over the $8$--$12$~dB decoding transition (onset $8.0$~dB), so SNR is highly informative; the learned JSCC feature F1 is \emph{essentially flat} at $\approx 0.89$ across $0$--$20$~dB, so the estimated SNR says almost nothing about whether a feature message will help on a given frame.

This flips the comparison (Table~\ref{tab:two_regime}). Both codecs are evaluated on the same $2{,}170$ test frames under the identical $200$-seed protocol, separating (a) whether the graceful-channel gain exists and resides in the cues---estimated leakage-free \emph{in-distribution} by $5$-fold cross-validation, every frame scored by a selector that never saw it and $\tau$ tuned on the training folds only---from (b) whether a \emph{deployed}, validate-frozen selector realises it, under the headline cross-split protocol.

\emph{(a) In-distribution.} Under LDPC + QAM the AWGN edge over the best SNR threshold is only
$+0.002$ F1 ($95\%$ CI $[+0.001,+0.003]$): small, and---as the feature ablation shows---not
attributable to the perception cues alone
(Section~\ref{sec:ablation}).\footnote{This in-distribution per-frame edge ($+0.002$) is
distinct from two other LDPC-regime quantities: the feature-ablation gain of the cues (Full
vs.\ channel-state-only, $+0.0090$, $95\%$ CI $[+0.0089,+0.0091]$, Section~\ref{sec:ablation}) and
the matched-payload margin over the retuned threshold ($+0.0032$, Section~\ref{sec:headline}). The
three measure the selector along the cue, the equal-bandwidth, and the per-frame axis respectively.
Under the corrected single-collaborator convention all three are positive, with the ordering
$+0.002 < +0.0032 < +0.0090$: the cue axis, smallest and indistinguishable from zero under the
retired full-collaborator accounting ($-0.0002$), is now the largest of the three --- weakening the
single message widened the gap the cues have to close, and they close part of it.} Under JSCC the best SNR threshold
collapses to the Fixed-$L$ operating point---it cannot gate on an uninformative SNR---and the
cue-based edge, measured by the in-distribution $k$-fold estimator, widens to $+0.022$ F1 ($[+0.020,+0.025]$) under AWGN, $+0.018$
($[+0.016,+0.020]$) under Rayleigh, and $+0.020$ ($[+0.018,+0.022]$) under OFDM on the test
split, and to $+0.012$--$+0.015$ on the Culver-City domain shift (all CIs exclude zero). In the separate $200$-realisation held-out comparison on validate frames, the
cue-based selector recovers $56$--$62\%$ of the clairvoyant oracle's headroom over Fixed $L$
across the three channels. The two quantities are distinct and are reported separately: under
AWGN the oracle headroom is $0.0291$ F1 and the selector recovers $+0.0181$ of it, with
$0.0275$ and $+0.0158$ under Rayleigh and $0.0281$ and $+0.0158$ under OFDM. In-distribution,
then, the graceful-channel decision is content-bound and carried by the ego-side cues---a
gain no SNR threshold reaches on these frames, and one that (b) below shows does not
survive the cross-split move.

\emph{(b) Cross-split deployment.} The signal does \emph{not} transfer: deployed on test the frozen selector over-selects the feature action (JSCC C-request rate $0.14\!\to\!0.42$) and falls just below Fixed-$L$---edge $-0.004$ (AWGN), $-0.008$ (Rayleigh), $-0.005$ (OFDM) on test, ${\approx}0$ on Culver-City. The gain is present but not \emph{portable}; capturing it under cross-domain deployment, e.g.\ by online or self-supervised adaptation~\cite{liu2024ssaw}, is future work.

\emph{Whether task-oriented cues are needed is determined by the feature codec's channel response.} Under a cliff codec the decision is channel-bound and SNR suffices; under a graceful codec it becomes content-bound and the ego-side cues---which predict whether object-level detection is already adequate---become the operative signal, in-distribution. Closing the cross-split transfer gap (b) is the remaining challenge.

\begin{table}[t]
\centering
\footnotesize
\setlength{\tabcolsep}{3pt}
\caption{Learned-selector edge over the best SNR-threshold rule on the same $2{,}170$ test
frames ($200$-seed protocol), reported two ways. \emph{In-dist.\ edge}: leakage-free $5$-fold
cross-validation on the evaluation split (each frame scored by a selector that never saw it,
$\tau$ tuned on the training folds only). \emph{Deployed edge}: selector trained on validate and
frozen, applied cross-split (the headline protocol). In-distribution the JSCC edge is
$9$--$11{\times}$ the LDPC-cliff edge---an SNR threshold cannot reach it---while under LDPC the
margin is the selector's function form, not a cue gain (cues add no significant F1,
Section~\ref{sec:ablation}). The deployed JSCC edge is negative: the validate-trained selector
does not transfer the graceful-channel gain across the domain shift
(Appendix~\ref{sec:jscc_aware}(b)).}
\label{tab:two_regime}
\resizebox{\columnwidth}{!}{%
\begin{tabular}{llcc}
\toprule
Feature codec & Channel & In-dist.\ edge & Deployed edge \\
\midrule
LDPC + QAM (cliff)      & AWGN     & $+0.002$~$[+0.001,+0.003]$ & $-0.002$~$[-0.003,-0.001]$ \\
Importance-map JSCC     & AWGN     & $\mathbf{+0.022}$~$[+0.020,+0.025]$ & $-0.004$~$[-0.007,-0.001]$ \\
Importance-map JSCC     & Rayleigh & $\mathbf{+0.018}$~$[+0.016,+0.020]$ & $-0.008$~$[-0.011,-0.006]$ \\
Importance-map JSCC     & OFDM     & $\mathbf{+0.020}$~$[+0.018,+0.022]$ & $-0.005$~$[-0.008,-0.003]$ \\
\bottomrule
\end{tabular}}
\end{table}

\section{Boundary of the Method: a Second Detection Backbone}\label{sec:second_backbone}

\noindent\emph{Convention.} This arm was trained and evaluated on the \textbf{full-collaborator} caches, i.e.\ before the corrigendum of Change-log P0: it is internally consistent (its own budget, its own $N_{cw}$, its own references) and its \emph{relative} conclusion stands, but its absolute F1 and payload figures are on a different accounting basis from the single-collaborator main experiment and are not tabulated beside it.

\noindent\emph{Positioning.} This appendix reports a \textbf{limitation}, not a generality result. It is placed here because its conclusion is negative: the selection procedure developed in the main text is in-sample effective on a second backbone but does not transfer to held-out splits under the equal-budget protocol. No claim in the main text depends on it.

Every result above uses one perception backbone. To test whether the selection
\emph{procedure}---not the particular features---is what carries the behaviour, we repeat the whole
pipeline on a second detector, SECOND~\cite{yan2018second}, under an \emph{equal-budget controlled}
comparison: the feature message is charged the same per-frame channel budget as the mainline
($B_F\equiv0.99$~Msym, the same $N_{\mathrm{cw}}$ and the same BLER tables), so the two backbones are
compared at equal channel cost. We do \emph{not} claim that SECOND's feature tensor compresses to
that size, and we declare no codec for it.\footnote{The transmitted tensor was measured, not assumed:
SECOND's cross-link BEV tensor is $6{,}758{,}400$ elements per collaborator before compression and
$352{,}000$ at the auto-encoder bottleneck. Both are recorded as measurements; neither is used as
the operative payload under the equal-budget convention.}

Both SECOND checkpoints are the official OpenCOOD ones, and we verify them before use rather than
trusting the file: after a lossless axis-order conversion required by the spconv version in our
environment, official inference reproduces the published AP@$0.7$ to $+0.0002$ on both splits for the
late-fusion model ($0.7752$ vs.\ $0.775$; $0.6822$ vs.\ $0.682$) and to within $+0.0019$ for the
intermediate model. The three per-frame branches ($E$, $L$, $F$) are then re-derived with the same
scorer, the same canonical union ground truth and the same IoU-$0.5$ convention as the mainline, and
the entire downstream pipeline---grid expansion, scene-level $9$-fold LOSO, the budget walk, and the
$200$-realisation paired replay---is re-run with the mainline code, each stage first verified to
reproduce its own committed mainline product bit-for-bit.

The result is that \method{} is \textbf{in-sample effective but does not transfer under the
equal-budget protocol}. Table~\ref{tab:second_backbone} gives the paired comparison against an
SNR-threshold rule tuned on validate for the same target budget. On validate---the split the selector is trained on---it is ahead
at every budget ($+0.0058$ to $+0.0095$ F1). On both held-out splits it is behind at every budget,
with all confidence intervals excluding zero.

\begin{table}[t]
\centering
\footnotesize
\setlength{\tabcolsep}{3pt}
\caption{Second backbone (SECOND), equal-budget controlled. Paired $\Delta$F1 of the frozen selector
against the SNR-threshold rule tuned to the \emph{same} budget, over the same $200$ CSI samplings;
$95\%$ frame-level paired bootstrap intervals. Positive favours \method{}. \emph{Second-backbone
arm, not deployed}; descriptive, no decision is taken from this table.}
\label{tab:second_backbone}
\begin{tabular}{lccc}
\toprule
Split & $B_{\max}=0.10$ & $B_{\max}=0.20$ & $B_{\max}=0.30$ \\
\midrule
Validate (in-sample) & $+0.0058$ & $+0.0095$ & $+0.0081$ \\
                     & \tiny$[+0.0057,+0.0059]$ & \tiny$[+0.0094,+0.0096]$ & \tiny$[+0.0081,+0.0082]$ \\
OPV2V test           & $-0.0014$ & $-0.0047$ & $-0.0066$ \\
                     & \tiny$[-0.0015,-0.0013]$ & \tiny$[-0.0048,-0.0046]$ & \tiny$[-0.0067,-0.0065]$ \\
Culver-City          & $-0.0029$ & $-0.0117$ & $-0.0167$ \\
                     & \tiny$[-0.0030,-0.0028]$ & \tiny$[-0.0119,-0.0115]$ & \tiny$[-0.0169,-0.0165]$ \\
\bottomrule
\end{tabular}
\end{table}

Two mechanisms account for the gap, and both are visible in the arm's own numbers. First, the
\textbf{$L$--$F$ margin is narrower on this backbone}: SECOND's object-level branch scores $0.869$ /
$0.879$ / $0.879$ mean per-frame F1 on validate / test / Culver-City against a compressed-feature
branch at $0.897$ / $0.904$ / $0.939$, so there is simply less to gain from choosing correctly, and
the same selection error costs proportionally more. Second, the \textbf{budget walk selects a
different candidate family}: all three frozen SECOND selectors carry \texttt{class\_weight=balanced},
whereas all three mainline selectors carry \texttt{class\_weight=None}. The reweighted family fits
the in-sample class mix well---agreement with the oracle is $0.833$ on validate---but transfers
poorly, falling to $0.578$ on test and $0.534$ on Culver-City, where the selector degenerates to
near-always-$L$ ($E$ recall $0.027$ on test and $0.000$ on Culver-City). The $E$-collapse of
Section~\ref{sec:pareto} is therefore \emph{worse} on this backbone, not milder: under Rayleigh the
selector's $\rho_E$ is $0.002$ (test) and $0.000$ (Culver-City) while the oracle spends $0.157$ and
$0.133$. The large payload reductions this arm shows ($67$--$90\%$) follow from that collapse to $L$
and are not evidence of good selection.

We report this as a boundary of the present method rather than a tuning target: nothing was
retrained and no threshold was adjusted after seeing these numbers.


\bibliographystyle{IEEEtran}
\bibliography{refs}

\end{document}
